\section{Faktorisierung}

\begin{df}\label{skript:2.1}
	Sei $ f \in \H^2 $.
	Wir bezeichnen $ f $ als \textit{innere Funktion}, falls $ | f | = 1 $ fast überall auf $ \mathbb{T} $ gilt.
	Wir nennen $ f $ eine \textit{äußere Funktion}, falls
	\begin{align*}
	E_f := \bigvee_{n \in \N_0} \mathscr{S}^n f = \H^2
	\end{align*}
	gilt.
\end{df}

\begin{genericthm}{Faktorisierungstheorem}\label{skript:2.2}
	Sei $ f \in \H^2 \setminus \lbrace 0 \rbrace $.
	Dann lässt sich $ f $ bis auf eine multiplikative Konstante eindeutig als Produkt einer inneren Funktion $ f^i $ und einer äußeren Funktion $ f^a $ beschreiben.
	D.h. es gilt $ f = f^i f^a $ und $ E_f = f^i \H^2 $.
\end{genericthm}

\begin{proof}
	Wegen $ f \neq 0 $ erhalten wir  mit \ref{skript:1.3}
	\begin{align*}
	E_f = \Theta \H^2	
	\end{align*}
	mit $ |\Theta| = 1  $ f. ü. auf $ \mathbb{T} $.
	Damit ist $ \Theta  $ eine innere Funktion. Wir setzen $ f^i := \Theta $ und $ f^a := f \overline{\Theta} $, womit gilt $ f = f^i f^a $. 
	Wir müssen noch nachweisen, dass $ f^a $ eine äußere Funktion ist.
	Da $ f^i $ als Multiplikationsoperator eine Isometrie auf $ \L^2 $ bzw. $ \H^2 $ ist, erhalten wir 
	\begin{align*}
	E_f
	= \bigvee_{n \in \N_0} z^n f^i f^a 
	= f^i \bigvee_{n \in \N_0} z^n f^a 
	= f^i E_{f^a}.
	\end{align*}
	Damit folgt $ f^i H^2 = f^i E_{f^e} $ und damit auch $ H^2 = E_{f^e} $. Also ist $ f^e $ eine äußere Funktion.
	Die Eindeutigkeit erhalten wir aus Beurling-Helson und \ref{skript:1.2}.
\end{proof}

\begin{genericthm}{Folgerung}\label{skript:2.3}
	Sei $ f \in \H^2$.
	Dann gilt:
	\begin{align*}
	f \ \text{ist äußere Funktion} \
	\Leftrightarrow
	\left(g \in \H^2, \frac{g}{f} \in \L^2 \Rightarrow \frac{g}{f} \in \H^2 \right)
	\end{align*}
\end{genericthm}

\begin{proof}
	Zunächst beginnen wir mit der Hinrichtung.
	Sei $ f $ eine äußere Funktion, $ g \in \H^2$ und $ \frac{g}{f} \in \L^2 $. Wegen $ E_f = \H^2 $ existiert eine Folge $ p_n $ von Polynomen, sodass
	\begin{align*}
	\lim\limits_{n \to \infty} \| p_n f - 1 \|_2 = 0
	\end{align*}
	gilt. Damit erhalten wir 
	\begin{align*}
	\left\| p_n g - \frac{g}{f} \right\|_1 
	= 
	\left\| (p_n f - 1 ) \frac{g}{f} \right\|_1 
	\leq
	\| p_n f - 1 \|_2 \left\| \frac{g}{f} \right\|_2
	\stackrel{n \to \infty}{\rightarrow} 0
	\end{align*}
	mit der Hölderungleichung. Damit gilt
	\begin{align*}
	\left(\frac{g}{f}\right)^{\hat{}}(n) = 0
	\end{align*}
	für $ n <0 $. Somit gilt $ \frac{g}{f} \in \H^2$.\\
	Nun kommen wir zur Rückrichtung.
	Wir setzen $ g = f^a $.
	Dann gilt
	\begin{align*}
	\frac{g}{f} = \frac{1}{f^i} = \overline{f^i} \in \L^2
	\end{align*}
	und mit der Voraussetzung folgt $ \overline{f^i} \in \H^2 $.
	Wegen $ f^i \in \H^2 $ folgt, dass $ f^i $ konstant ist.
	Somit ist $ f $ eine äußere Funktion.
\end{proof}

\begin{genericthm}{Folgerung}\label{skript:2.4}
	\begin{enumerate}[label=\text{\textbf{(\arabic*)}}]
		\item
		Seien $ f$ und $g $ äußere Funktionen mit $ |f| = |g| $ fast überall.
		Dann folgt $ f = \lambda g $ mit $ \lambda \in \mathbb{T} $.
		
		\item Sei $ f \in \L^2 $ und $ \mathscr{S}^{-1} E_f \neq  E_f$.
		Dann existiert eine äußere Funktion $ g $ mit $ |f | = |g|  $ fast überall auf $ \mathbb{T} $.
		
%		\item
%		Falls für $ h \in \L^2 $ mit $ h >0 $ 
%		\begin{align*}
%		\int \limits_{\mathbb{T} } \log h > - \infty
%		\end{align*}
%		gilt, folgt $ |h| = |f| $ fast überall für eine äußere Funktion $ f $.
	\end{enumerate}
\end{genericthm}

\begin{proof}
	\begin{enumerate}[label=\text{\textbf{(\arabic*)}}]
		\item 
		Seien $ f $ und $ g $ äußere Funktionen mit $ |f| = |g| $ fast überall.
		Dann ist mit \ref{skript:2.3}
		\begin{align*}
		\frac{g}{f} \in \H^2 \ \wedge \overline{ \left(\frac{f	}{g} \right)} \in \H^2,
		\end{align*}
		womit folgt, dass $ \frac{g}{f} $ konstant ist. Somit gilt die Aussage.
		\item
		Wir wenden Beurling-Helson an. Damit existiert ein $ \Theta \in \L^2 $ mit $ |\Theta| = 1 $ fast überall, sodass
		\begin{align*}
		E_f = \Theta \H^2
		\ \Leftrightarrow \ 
		\overline{\Theta} E_f = E_{\overline{\Theta} f} = \H^2
		\end{align*}
		gilt. Damit wissen wir, dass ein $ g \in \H^2 $ mit
		\begin{align*}
		\iota(g) = \overline{\Theta} f
		\end{align*}
		existiert. Die Einbettung $ \iota $ ist eine Isometrie, wodurch die Aussage direkt mit \ref{skript:2.2} folgt.
	\end{enumerate}
\end{proof}

\begin{genericthm}{Folgerung}\label{skript:2.5}
	Sei $ f \in \L^2 $. Dann gelten:
	\begin{enumerate}[label=\text{\textbf{(\arabic*)}}]
		\item 
		$ \mathscr{S} E_f \neq E_f \ \Leftrightarrow  \ \log|f|  \in \L^1$
		\item 
		$ E_f = \L^2   \ \Leftrightarrow \ \log|f| \notin \L^1$ mit $ f \neq 0 $ f.ü.
	\end{enumerate}
\end{genericthm} 

\begin{proof}
	\begin{enumerate}[label=\text{\textbf{(\arabic*)}}]
		\item 
		Sei $ \mathscr{S} E_f \neq E_f $. Mit dem selben Argument wie im zweiten Teil der letzten Folgerung erhalten wir
		\begin{align*}
		\H^2 = E_{ \overline{\Theta} f}
		\end{align*}
		mit $ \overline{\Theta} f \in \iota(\H^2) $.
		Dies ist äquivalent zu $ \log| \overline{\Theta} f | = \log |f| \in \L^1 
		$.
		\item 
		Wegen $ E_f = \L^2 $ gilt $ \mathscr{S} E_f = E_f $, womit direkt $ \log|f| \notin \L^1 $ folgt.  
		Wegen 
		\begin{align*}
		E_f \subset \chi_e \L^2, \quad 
		e = \lbrace \xi \in \mathbb{T}\ | \ f(\xi) \neq 0 \rbrace
		\end{align*}
		gilt $ f\neq 0 $ fast überall.
		Anderseits gilt $ \log |f| \notin \L^1 $ mit der ersten Aussage $ \mathscr{S} E_f = E_f $. Beurling-Helson liefert dann $ E_f = \chi_e \L^2 $.
		Mit $ \lambda_N(e) = 1 $ folgt dann $ E_f = \L^2 $.
 	\end{enumerate}
\end{proof}

\begin{df}\label{skript:2.6}
	Eine innere Funktion $ \Theta_1 $ \textit{teilt} die innere Funktion $ \Theta_2 $, falls
	\begin{align*}
	\frac{\Theta_2}{\Theta_1} \in \H^2
		\end{align*}
		gilt. Damit ist $ \frac{\Theta_2}{\Theta_1} $ eine innere Funktion.
		Wir nennen $ \Theta $ das kleinste gemeinsame Vielfache ($ \kgv $) von $ \Theta_1 $ und $ \Theta_2 $, falls $ \Theta $ von $ \Theta_1 $ und $ \Theta_2 $ geteilt wird und jede innere Funktion $ \Theta^\prime $, welche von $ \Theta_1 $ und $ \Theta_2 $ geteilt, auch von $ \Theta $ geteilt wird.
		
\end{df}

\begin{genericthm}{Satz}\label{skript:2.7}
	
	\begin{enumerate}[label=\text{\textbf{(\arabic*)}}]
		\item Es gilt $ \Theta_2 \H^2 \subset \Theta_1 \H^2  $ genau dann, wenn
		$ \Theta_1$ teilt $ \Theta_2 $ gilt.
		
		\item
		Es gilt
		\begin{align*}
		\Theta_1 \H^2 \cap \Theta_2 \H^2 = \Theta \H^2,
		\end{align*}
		wobei $ \theta $ das kleinste gemeinsame Vielfache von $ \Theta_1 $ und $ \Theta_2$ ist.
	\end{enumerate}
	
\end{genericthm}

\begin{proof}
	\begin{enumerate}[label=\text{\textbf{(\arabic*)}}]
		\item
			Wegen $ \Theta_2 \H^2 \subset \Theta_1 \H^2 $ existiert eine innere Funktion $ f $ mit $ \Theta_2 = \Theta_1 f $.\\
			Aufgrund von
			\begin{align*}
			\Theta_2 \H^2 = \Theta_1 f \H^2 \subset \Theta_1 \H^2
			\end{align*}
			erhalten wir $ \Theta_2 \in \Theta_1 \H^2 $. Somit gilt auch 
			$ \frac{\Theta_2}{\Theta_1} \in \H^2 $.
		\item 
			Der Raum $ \Theta_1 \H^2 \cap \Theta_2 \H^2 $ ist $ S $-invariant ist. Des weiteren ist $\Theta= \Theta_1 \Theta_2 \neq 0 $ in diesem Unterraum enthalten.
			Demnach gilt $\Theta_1 \H^2 \cap \Theta_2 \H^2 = \Theta \H^2  $
			Es folgt direkt, dass $ \Theta $ von $ \Theta_1 $ und $ \Theta_2 $ geteilt wird.
			Wenn nun ein $ \Theta^\prime $ von $ \Theta_1$ und $ \Theta_2 $ teilbar ist, erhalten wir mit
			\begin{align*}
			\Theta^\prime \H^2 \subset \Theta_1 \H^2 \ \wedge \ \Theta^\prime \H^2 \subset \Theta_2 \H^2
			\end{align*}
			sofort $ \Theta^\prime \H^2 \subset \Theta \H^2 $.
			Damit ist $ \Theta  $ das kleinste gemeinsame Vielfache von $ \Theta_1 $ und $ \Theta_2 $.
	\end{enumerate}
\end{proof}

 
